\chapter{Task identification}

In this chapter we describe the phases that led to the acquisition of the domain
knowledge, the identification of the requirements to be implemented, and the
definition of the reasoning strategy.

The task of BeerEX is, given a user, to recommend the beer style(s) that best
suit that user in the current situation. The recommendation is grounded on three
layers of information, acquired through a guided dialogue: the \emph{scenario} of
the user, their personal \emph{preferences}, and the \emph{meal} they are about
to enjoy. From the answers, the expert rules infer a desired beer profile, and a
selection stage ranks the candidate styles by how well they match it, reasoning
under uncertainty through one of two interchangeable calculi.


\section{Knowledge Acquisition}

The acquisition of the knowledge of the application domain was carried out by
drawing on two complementary kinds of source:

\begin{itemize}
    \item \emph{Authoritative guides}: the beer-style and food-pairing material
    published by CraftBeer.com on behalf of the Brewers
    Association~\cite{craftbeer-styles,craftbeer-food}, namely the \emph{Beer Styles
    Study Guide} (the descriptive and numeric characterisation of each style:
    colour, alcohol, bitterness, fermentation, flavour) and the food-pairing
    tables (the recommended style for each food).

    \item \emph{Domain heuristics}: the rules of thumb that relate the user's
    scenario---age, drinking habit, smoking habit, season---to the strength
    and flavour profile of an enjoyable beer, distilled into the \emph{why}
    rationales attached to the corresponding questions.
\end{itemize}

Unlike the descriptive guides, which characterise the styles, the recommendation
task requires eliciting the \emph{user's} situation. The system therefore
acquires the knowledge needed to advise the user through a dialogue organised in
three layers:

\begin{enumerate}[nosep]
    \item the \textbf{user scenario}---age, whether a regular beer drinker,
    whether smoking, and the season---factors that shift, weakly, the expected
    strength and flavour of an enjoyable beer;
    \item the \textbf{personal preferences}---preferred colour, alcohol level,
    carbonation and flavour;
    \item the \textbf{meal}, when the user is about to eat---main component, type
    and cooking method, refined through increasingly specific questions down to a
    single dish.
\end{enumerate}

Each answer asserts a fact in working memory, from which the expert rules infer
\emph{evidence} for the desired beer profile. The questions handled by the
system, grouped by layer, together with the attribute each elicits and the set
of admissible answers, are classified in Table~\ref{tab:questions}.

\begin{longtable}{p{0.27\linewidth} p{0.27\linewidth} p{0.38\linewidth}}
\caption{Classification of the question groups handled by the expert system, with
the attribute each elicits and its admissible answers.}\label{tab:questions}\\
\toprule
\textbf{Layer / group} & \textbf{Attribute elicited} & \textbf{Admissible
answers} \\
\midrule
\endfirsthead
\toprule
\textbf{Layer / group} & \textbf{Attribute elicited} & \textbf{Admissible
answers} \\
\midrule
\endhead
\midrule
\multicolumn{3}{r}{\footnotesize\itshape continued on next page}\\
\endfoot
\bottomrule
\endlastfoot

\multicolumn{3}{l}{\itshape Scenario (asked in random order)}\\
\addlinespace
Age & \code{which-age} & 18--23, 24--29, $\geq$30. \\
Regular drinker & \code{regular-beer-drinker} & yes, no. \\
Smoking habit & \code{smoker} & yes, no. \\
Season & \code{which-season} & autumn, spring, summer, winter. \\
\addlinespace
\midrule
\multicolumn{3}{l}{\itshape Preferences (asked in random order)}\\
\addlinespace
Colour & \code{preferred-color} & pale, amber, brown, dark, don't mind. \\
Alcohol & \code{preferred-alcohol} & low, mild, high, very high, don't mind. \\
Carbonation & \code{preferred-carbonation} & low, medium, high, don't mind. \\
Flavour & \code{preferred-flavor} & clean, sweet, bitter, roasty, fruity,
spicy, sour, don't mind. \\
\addlinespace
\midrule
\multicolumn{3}{l}{\itshape Meal (depth-first chain)}\\
\addlinespace
Main meal & \code{main-meal} & pizza, entrée, cheese, dessert, other. \\
\addlinespace
Pizza topping & \code{pizza-topping} & classic, meat, vegetables, cheese,
other. \\
Spicy meat & \code{meat-topping-is-spicy} & yes, no. \\
Roasted vegetables & \code{vegetables-topping-are-roasted} & yes, no. \\
\addlinespace
Entrée component & \code{which-entrée} & grain, legumes, fish, meat,
vegetables, fats, other. \\
Fish type & \code{which-fish} & shellfish, other. \\
Mild shellfish & \code{shellfish-is-mild} & yes, no. \\
Shellfish type & \code{which-shellfish} & mussels, oysters, other. \\
Meat type & \code{which-meat} & rich, poultry, game, other. \\
Meat cooking method & \code{meat-cooking-method} & braised, grilled, roasted, other. \\
Rich meat & \code{which-rich} & beef, lamb, pork, other. \\
Beef type & \code{which-beef} & bresaola, other. \\
Pork cut & \code{which-pork} & loin, prosciutto, speck, mortadella, sausage, other. \\
Sausage type & \code{which-sausage} & capocollo, soppressata, salame piccante,
other. \\
Poultry type & \code{which-poultry} & chicken, other. \\
Game type & \code{which-game} & wild, birds. \\
Vegetables type & \code{which-vegetables} & root, other. \\
Vegetables cooking method & \code{vegetables-cooking-method} & grilled, other. \\
Other vegetables & \code{which-other-vegetables} & mushrooms, other. \\
Fats type & \code{which-fats} & vegetable, other. \\
\addlinespace
Cheese style & \code{which-cheese-style} & fresh, semi-soft, firm/hard, blue,
natural rind, washed rind, don't know. \\
Fresh cheese & \code{which-fresh-cheese} & Mascarpone, Chèvre, Feta,
Cream Cheese, other. \\
Semi-soft cheese & \code{which-semi-soft-cheese} & Mozzarella, Colby, Havarti,
Monterey Jack, other. \\
Firm/hard cheese & \code{which-firm/hard-cheese} & Gouda, Cheddar, Swiss,
Parmesan, other. \\
Gouda type & \code{which-type-of-Gouda} & aged, smoked, other. \\
Cheddar colour & \code{which-color-of-Cheddar} & white, other. \\
Cheddar seasoning & \code{which-Cheddar-seasoning} & mild, medium, aged,
other. \\
Swiss type & \code{which-type-of-Swiss} & Emmental, Gruyère, other. \\
Blue cheese & \code{which-blue-cheese} & Stilton, other. \\
Natural-rind cheese & \code{which-natural-rind-cheese} & Brie, Camembert,
Triple Crème, Mimolette, Stilton, other. \\
Washed-rind cheese & \code{which-washed-rind-cheese} & Taleggio, other. \\
\addlinespace
Dessert type & \code{which-dessert} & creamy, fruit, chocolate, other. \\
Fruity creamy dessert & \code{creamy-dessert-with-fruit} & yes, no. \\
Chocolate type & \code{which-chocolate} & white, milk, semisweet, bittersweet,
unsweetened/bitter, don't know. \\
\end{longtable}


\section{Requirements}

In the following sections we provide a brief description of the requirements that
the system must satisfy and of the features to be implemented.


\subsubsection{Programming Language}

The expert system is implemented in CLIPS (the \emph{C Language Integrated
Production System}), a forward-chaining production-system tool well suited to
rule-based expert systems.

By using CLIPS as the programming language, the knowledge about the identified
objects and relations can be expressed declaratively, in the form of facts and
production rules, while the questions to be posed to the user---needed to derive
further information useful to the inference process---are themselves encoded as
rules that assert the dialogue state. The inference engine matches the rules
against working memory by means of the \emph{Rete} algorithm, which makes the
pattern matching efficient even as the number of facts and rules grows.

We rely on a stock, modern CLIPS (the 6.4.x line), without forking the engine,
so that the knowledge base remains portable. For the deployment as a Telegram
bot, the engine is driven from Python through \code{clipspy}, the official CLIPS
bindings, which embed the very same engine used by the command-line interpreter,
so that the behaviour of the bot and of the CLI coincide.

Two behavioural features required by the system deserve a dedicated discussion:
the handling of uncertainty and an explanation module. They are discussed below.


\subsubsection{Handling of Uncertainty}

The first behavioural feature concerns the handling of uncertainty.

A categorical kind of reasoning, in which every element can only be either true
or false, turns out to be too restrictive for the task at hand. The suitability
of a beer is a matter of degree: a given scenario, preference or dish does not
single out one style, but rather makes a whole profile---a colour, an alcohol
level, a flavour---\emph{more or less} desirable. The evidence supporting the
desired profile is therefore inherently graded, and the system must be able to
accumulate it and to weigh it against the measured characteristics of each
candidate style.

From these premises, the need emerges to move away from the categorical nature of
assertions and to attach a degree of confidence to every piece of evidence and to
every recommendation derived from it.

BeerEX provides two interchangeable calculi for reasoning with this graded
evidence, selectable through a single backend flag. The first is the
\emph{certainty-factor} (CF) calculus of MYCIN, in which each piece of evidence
carries a certainty factor in $[-1,1]$ and contributions to the same conclusion
are merged with the MYCIN combination function. The CF weights, however, are
assigned by the expert and the source guides do \emph{not} quantify them; this
intrinsic arbitrariness is precisely what motivates the second calculus, a
\emph{fuzzy-logic} backend whose membership functions are derived from data---by
crossing the symbolic labels of the dataset with the numeric ranges of the
guide---so that the matching of the candidate styles against the desired profile
rests on objective, data-grounded quantities rather than on hand-assigned
weights.

The two backends are described in depth in the chapter devoted to knowledge
conceptualisation; here it suffices to note that both consume the same evidence
and differ only in how they combine it and match it to the styles.


\subsubsection{Explanation Module}

The second behavioural feature concerns the implementation of an explanation
module.

In an expert system, the explanation capability is of fundamental importance.
Users tend not to trust an artificial system slavishly, nor to accept blindly a
recommendation produced in an uncontrollable way. The capability to justify both
the questions and the result increases the trust the user places in the advice of
the system and enables them to discover a possible flaw in the reasoning.

In keeping with this requirement, the expert system provides explanations on two
distinct occasions. During the dialogue, at every question the user may ask for
\emph{help}---an explanation of the question, sometimes with illustrative
images---or for \emph{why} the question is being asked, so as to suit different
levels of expertise. At the end of the consultation, the system returns a
twofold explanation of the result: a narrative that recounts how the user
scenario, the preferences and the meal shaped the desired profile, and---in fuzzy
mode---a per-beer rationale that lists the factors that matched each
recommendation and how strongly.


\section{Reasoning Strategy}

As far as the type of inference is concerned, the choice fell on forward
chaining, the inference natively provided by CLIPS: starting from the known facts
contained in working memory, the production rules are applied forward in order to
deduce new facts, until no further rule is applicable.

The dialogue drives the inference in three layers. The \emph{scenario} and
\emph{preference} questions are independent of one another and are therefore
asked in random order, by setting the conflict-resolution strategy to
\code{random} while their rules are eligible to fire. The \emph{meal} questions,
on the contrary, form a depth-first chain: each answer enables the next, more
specific question, narrowing the dialogue from the main course down to a single
dish---so the strategy is switched to \code{depth} once the meal phase begins.

Each answer asserts a fact in working memory through the \code{relation-asserted}
of the question. The knowledge rules, which sit at a medium-low salience, then
fire on these facts and assert the \emph{evidence} for the desired profile, in
the form \code{(attribute (name best-\textit{X}) (value \textit{v}) (certainty
\textit{w}))}, for the colour, alcohol, carbonation, flavour and fermentation
axes, together with direct suggestions of a specific style or beer name and the
narrative fragments of the explanation.

Because the dialogue allows the user to go \emph{back} and revise a previous
answer, the working memory must be kept consistent with the answers actually
given. When the user goes back, the dialogue state machine retracts the fact
asserted by the revised question and refreshes the affected rules, so that
everything inferred from the superseded answer is removed and re-derived from the
new one; in random mode the question rules are refreshed so they may be asked
again, and if the consultation had already reached its final state the selection
rules are refreshed as well.

Finally, once all the evidence has been gathered, a selection stage at low
salience ranks the styles. Contributions to the same \code{best-*} value are
merged with the certainty-factor combination function; each candidate style is
then scored against the desired profile---by accumulating the certainties of the
matching axes in CF mode, or by a data-grounded weighted-mean satisfaction in
fuzzy mode---and the most appropriate styles are returned, each with its
confidence and a serving hint, followed by the explanation.

The reasoning strategy introduced here is taken up again and documented in depth
in the chapters devoted to the description of the inference engine.
